% Root and FGH independent proof/code audits passed on 2026-11-15.
\begin{statement}{C.37}{Proved}{The remaining primes five, seven and eleven}
Let $Q$ be the quotient of the full gentle group by the
closed normal subgroup generated by its absolute Galois image.
Then
\[
\begin{array}{c|rrr}
p&5&7&11\\ \hline
\dim_{\mathbb F_p}\operatorname{Hom}_{\mathrm{cts}}
             (Q,\mathbb F_p)&\geq1&\geq1&\geq513 .
\end{array}                                                     \tag{C37.1}
\]
Every displayed independent family has its independent
values witnessed on global multiplier one.
Together with C.34--C.36, this gives a nonzero continuous
arithmetic-trivial character of the full gentle group
to $\mathbb F_p$ for every prime $p\geq5$.
The general odd-primary families in this statement are
not asserted to vanish on the pentagon group.
\end{statement}
\begin{proof}
We allow repeated tame eigencharacters, which were avoided
in C.34 by truncating at degree $p-1$.
For the three primes choose
\[
                 (p,m)=(5,12),(7,12),(11,20).
                                                               \tag{C37.2}
\]
Use the actual local image $\mathcal G$, its
multiplier-one subgroup $\mathcal U$, and its conjugator
filtration from C.33, now at the indicated prime.
The completed P.2 correction and P.4 lifting give
$\mathcal U_n/\mathcal U_{n+1}=S_n(\mathbb Z_p)$
in every degree used here.
Set $U=\mathcal U/\mathcal U_{m+1}$ and
$G=\mathcal G/\mathcal U_{m+1}$.
Their integral Magnus operators raise degree by at
least three, so $U$ has nilpotency class at most
$\lfloor m/3\rfloor$, respectively $4,4,6$.
Each is strictly less than $p$, and the actual
Lazard Lie lattice $\mathfrak u=\log U$ is integral.

Let $t$ be a tame order-$(p-1)$ element in the local
Galois image, with primitive multiplier $\omega$.
Its averaging projector onto eigenvalue one is
\[
                 e=\frac1{p-1}
                       \sum_{j=0}^{p-2}\operatorname{Ad}(t)^j .
\]
Although the same eigencharacter can occur in several
degrees, we can choose an integral eigenbasis adapted
to the filtration.  Indeed averaging splits each
eigenspace, and each successive filtered quotient
is a free $\mathbb Z_p$ module; lift a basis of the
quotient and project to its required eigencharacter.
Thus each basis vector has a specified first degree
$n$, of multiplicity $a_n$, and eigencharacter
$\omega^n$.
It follows that
\[
 \operatorname{rank}e\mathfrak u
        =\sum_{\substack{2\leq n\leq m\\p-1\mid n}}a_n.
                                                               \tag{C37.3}
\]

The following count concerns actual filtered brackets,
not only their leading terms.
For two such basis vectors with first degrees $i,j$,
their bracket has eigencharacter $\omega^{i+j}$
and filtration at least $i+j$.  It can contribute to
$e[\mathfrak u,\mathfrak u]$ only if
$p-1\mid i+j\leq m$.
There is at most one resulting vector for each pair
of basis vectors.  Antisymmetry removes the diagonal
pair.  Hence
\[
 \operatorname{rank}e[\mathfrak u,\mathfrak u]
       \leq\sum_{\substack{2\leq n\leq m\\p-1\mid n}}b_n.
                                                               \tag{C37.4}
\]
Higher terms of a lower-degree bracket can change
that one vector, but cannot increase this count.
The Lie derived ideal is the linear span of all
such pair brackets, so this also counts every
possible hidden filtered relation.

Here are all the necessary dimensions:
\[
\begin{array}{c|c|r|r|r}
p&\text{invariant degrees with possible nonzero space}
 &\operatorname{rank}e\mathfrak u
 &\text{upper bound in (C37.4)}&\text{difference}\\ \hline
5&8,12&8&7&1\\
7&12&7&6&1\\
11&10,20&820&307&513 .
\end{array}                                                     \tag{C37.5}
\]
The omitted invariant degrees $4,6$ have zero space.
At degree twenty, the exact character formula of C.34
gives $a_{20}=818$ and $b_{20}=305$.
The program \texttt{pbc\_small\_prime\_characters.py}
checks these rows and the class bounds, retains the
input hashes, and independently recomputes the
lower-degree bracket ranks used next.

The lower invariant degrees cause no nontrivial
weight extension in the rational abelianization.
For $p=5$, the degree-eight space has dimension one
and is the span of the bracket of degrees three and
five.  For $p=7$ there is no nonzero lower invariant
degree.  For $p=11$, the degree-ten space has
dimension two and is spanned by the brackets of
degrees three and seven.
These are the exact bracket ranks in (C33.3).
Lift their factors to the appropriate tame eigenspaces
of the actual lattice.
The resulting actual brackets span each indicated
lower invariant graded piece over $\mathbb Q_p$.
For any invariant lattice vector, subtract a rational
combination of those actual brackets to cancel its
first lower component.  Continue if necessary.
There is no other lower invariant graded piece.
Consequently the natural map
\[
 \mathfrak u_m\otimes\mathbb Q_p
 \longrightarrow
      e\bigl(\mathfrak u/[\mathfrak u,\mathfrak u]\bigr)
                      \otimes\mathbb Q_p
                                                               \tag{C37.6}
\]
is onto.

Let $R$ be the integral lattice of functionals
$\ell:\mathfrak u\to\mathbb Z_p$ which annihilate
$[\mathfrak u,\mathfrak u]$ and are invariant under
$t$.  Equations (C37.3)--(C37.5) give ranks at least
$1,1,513$.
This is a saturated submodule of the integral dual
of $\mathfrak u$: if $p\ell$ satisfies both defining
linear conditions, so does $\ell$.
Thus an integral basis of $R$ remains linearly
independent after reduction modulo $p$.

Each such reduction is invariant under the whole $G$.
Conjugation on the Lie abelianization factors through
$G/U=\mathbb Z_p^\times$ and hence commutes with $t$.
By (C37.6), its invariant part over $\mathbb Q_p$
is generated by the top degree, where multiplier
$a$ acts exactly by $a^m$.
Therefore, for $\ell\in R$ and $g\in G$,
\[
          \ell(\operatorname{Ad}(g)v)=
                         \lambda_p(g)^m\ell(v).                 \tag{C37.7}
\]
This equality first holds over $\mathbb Q_p$ and
hence holds in the torsion-free integral target
$\mathbb Z_p$ as well.
In each row of (C37.2), $p-1$ divides $m$.
Reduction modulo $p$ makes the factor in (C37.7)
equal to one.
The reductions thus give the required number of
independent invariant ordinary homomorphisms
$U\to\mathbb F_p$.

The rest of the full-group argument is exactly the
already proved one in C.33: the continuous five-term
sequence and
$H^2_{\mathrm{cts}}(\mathbb Z_p^\times,\mathbb F_p)=0$
extend each character to $G$.
Its arithmetic restriction is unramified away $p$;
Kronecker--Weber removes it by a scalar multiple of
$c_p(\lambda)$ from (C34.7).
The adjustment vanishes on global multiplier one.
All elements of the actual local multiplier-one
image have full multiplier-one lifts by P.4, so
the independent values of the reductions remain
actual full-group witnesses.
This proves (C37.1) and the last assertion.
\end{proof}
